\chapter*{Introduction}
\addstarredchapter{Introduction}
%%%%%%%%%%%%%%%%%%%%%%%%%%%%%%%%%%%%%%%%%%%%%%%%%%%%%%%%%%%%%%%%%%%%%%%%

L'introduction précède le corps du mémoire. Cette partie essentielle permet
d'introduire le lecteur au sujet traité et de présenter les enjeux principaux.
Le texte se doit d'être accrocheur afin de susciter l'intérêt du lecteur pour
la suite. C'est un travail rédactionnel qui nécessite une attention particulière
pour produire un mémoire original et démontrer son implication dans la recherche.

Le mémoire est organisé comme suit :
Le premier chapitre présente le projet décisionnel dans son contexte général.
Le deuxième chapitre identifie les besoins d'affaires et expose l'architecture technique retenue.
Le troisième chapitre traite de la modélisation conceptuelle et physique de l'entrepôt de données.
Le quatrième chapitre détaille le processus d'intégration des données.
Le cinquième chapitre présente la conception, le développement et le déploiement de la plateforme décisionnelle.

\begin{tcolorbox}[
    enhanced,
    title={\textbf{Attention : Règles de ponctuation en français}},
    colback=yellow!10,
    colframe=red!75!black,
    colbacktitle=red!85!black,
    coltitle=white,
    fonttitle=\large\bfseries,
    boxrule=0.5mm
]

\textbf{Règles de ponctuation française :}

\begin{itemize}[leftmargin=*]
    \item Le point (.) : 
        \begin{itemize}
            \item Pas d'espace avant
            \item Un espace après
            \item Suivi d'une majuscule pour nouvelle phrase
        \end{itemize}
    
    \item La virgule (,) :
        \begin{itemize}
            \item Pas d'espace avant
            \item Un espace après
            \item Sépare les éléments d'une énumération ou les propositions
        \end{itemize}
    
    \item Le point-virgule (;) :
        \begin{itemize}
            \item Une espace fine avant avec \verb|~|
            \item Un espace après
            \item Sépare deux propositions liées par le sens
        \end{itemize}
    
    \item Les deux-points (:) :
        \begin{itemize}
            \item Une espace fine avant avec \verb|~|
            \item Un espace après
            \item Introduisent une énumération, une citation ou une explication
        \end{itemize}
    
    \item Les points de suspension (...) :
        \begin{itemize}
            \item Pas d'espace avant
            \item Un espace après
            \item Exactement trois points
        \end{itemize}
    
    \item Les points d'exclamation (!) et d'interrogation (?) :
        \begin{itemize}
            \item Une espace fine avant avec \verb|~|
            \item Un espace après
            \item Pas de point après ces signes
        \end{itemize}
    
    \item Les guillemets français (« ») :
        \begin{itemize}
            \item Une espace fine après l'ouverture
            \item Une espace fine avant la fermeture
            \item Utilisés pour les citations directes
        \end{itemize}
    
    \item Les parenthèses () et les crochets [] :
        \begin{itemize}
            \item Pas d'espace à l'intérieur
            \item Collés au mot qu'ils encadrent
            \item Un espace avant l'ouverture et après la fermeture
        \end{itemize}
\end{itemize}

\textbf{Note importante :} En typographie française, on distingue l'espace normal de l'espace fine. L'espace fine est utilisée avant les signes de ponctuation doubles (; : ! ? »), tandis que l'espace normal est utilisé après tous les signes de ponctuation.

\end{tcolorbox}

